\section{Uniform comparison with an inert parent}

Let $\mathcal X$ be the class of finite newborn birth environments,
allowing finitely many inert core-neighbour stubs. The tag $u$ has one
leaf and one core neighbour, which is either a real parent or an inert
stub. Write
\[
 A_m(s;x)=(\mathcal K_s^m\mathbf1)(x),\qquad
 A_m^\infty(s)=A_m(s;x_\infty),
\]
where $x_\infty$ is the tag with its leaf and an inert parent stub.
All generations count only future core descendants of the tag.
The potential-genealogy construction in Section~\ref{model:section} gives
$\|\mathcal K_s\|_\infty\le1/s$ and
$A_m(s;x)\ge[2(s+2)]^{-m}$.

\begin{theorem}[Uniform boundary removal]\label{growth:minorization}
For every $s>0$, every $m\ge1$, and every $x\in\mathcal X$,
\begin{equation}\label{growth:comparison}
 \boxed{
 A_m(s;x)\ge\gamma_s A_m^\infty(s),
 \qquad \gamma_s=\frac1{12(s+3)}.
 }
\end{equation}
The constant is independent of both the initial environment and the
genealogical generation.
\end{theorem}

The point of the proof is to create a permanent no-return event using
only exterior randomness. Merely conditioning on the absence of actual
parent arrivals would generally bias the descendant process.

\subsection{An exterior process with a positive permanent success probability}

Consider a birth--death process $X$ on the nonnegative integers,
absorbed at zero, with births $x\to x+1$ at rate $x$ and deaths at
rate one while $x>0$. At every death-clock ring, attach an independent
uniform mark $U$. Kill the construction if $U\le1/X(t-)$; otherwise
carry out the death. A ring at $X=1$ always kills it. The generator of
the surviving process, acting on functions with value zero at the
cemetery state, is
\begin{equation}\label{growth:killed-generator}
 Gh(x)=x\bigl(h(x+1)-h(x)\bigr)
 +(1-1/x)\bigl(h(x-1)-h(x)\bigr)-h(x)/x.
\end{equation}

\begin{lemma}\label{growth:survival}
Starting from $X(0)=1$, the probability of never being killed is at least
$1/6$.
\end{lemma}
\begin{proof}
Use the bounded nonnegative function
\[
 h(0)=0,\qquad h(1)=\frac16,\qquad
 h(x)=1-\frac4{3x}\quad(x\ge2).
\]
Direct substitution gives
\[
 Gh(1)=0,\qquad Gh(2)=\frac7{36},\qquad
 Gh(x)=\frac{x-3}{3x(x+1)}\ge0\quad(x\ge3).
\]
The process is nonexplosive, being dominated by a linear pure-birth
process. Dynkin's formula, or localization followed by bounded
convergence, gives
$\mathbb E[h(X(t));\,\text{not killed by }t]\ge h(1)=1/6$.
Since $0\le h\le1$, survival to every finite $t$ has probability at least
$1/6$. Let $t\to\infty$.
\end{proof}

\subsection{Proof of the uniform comparison}

If the parent is already inert, the tagged component has exactly the
inert-parent law, so~\eqref{growth:comparison} is immediate.
Suppose the parent $p$ is real.

First prepare a parent leaf without changing the tag. Wait for the first
ring among the tag, its sole leaf, and its parent. These are three
independent rate-one clocks. Declare preparation successful only when
the first ring is the tag's ring and it chooses the parent, rather than
the leaf. Until that ring the tag remains in its newborn state, even
if the parent receives arrivals from elsewhere. If $T$ is this first
ring time and $E_0$ the successful outcome, then
\begin{equation}\label{growth:preparation}
 \mathbb E[e^{-sT}\mathbf1_{E_0}]
 =\int_0^\infty \tfrac12 e^{-(s+3)t}\,dt
 =\frac1{2(s+3)}.
\end{equation}
On $E_0$ no tagged child has yet been born, the tag still has its one
leaf and one parent, and the parent has gained a leaf.

Restart the clocks at $T$. Lower-couple the parent's leaf count $L_p$
with the exterior process $X$ of Lemma~\ref{growth:survival}, starting
from one. Use a common parent-leaf birth Poisson random measure:
marks below $X$ create births in $X$, and marks below $L_p$ create the
actual parent-leaf births. This measure is independent of all future
tag-descendant driving randomness. At \emph{every} parent core-clock
ring, decrease $X$ by one while it is positive; the actual $L_p$
decreases by at most one. All other arrivals at the parent only increase
$L_p$. Therefore, up to killing,
\[
 L_p(t)\ge X(t),\qquad D_p(t)\ge X(t).
\]

At a parent core-clock ring, choose the tag as target precisely when
$U\le1/D_p(t-)$. Otherwise choose uniformly among the remaining current
neighbours. This is an exact uniform-neighbour construction. The
reserved event $U\le1/X(t-)$ contains every possible parent-to-tag
arrival. The process $X$ uses the parent ring times but not these marks.

Let $E$ be the event that the exterior construction is never killed.
It depends only on the parent-leaf birth measure, the parent core-clock
times, and the reserved marks. In particular, conditional on the
preparation history, it is independent of the future driving randomness
inside the tag's descendant component. By the lemma, its conditional
probability is at least $1/6$.

On $E$, the parent never sends a leaf to the tag. The tag may still
send leaves to the parent, but such outward events have no effect on
its retained component and cannot create a return arrival on $E$.
Its entire descendant history therefore has the inert-parent law,
using the unchanged independent interior drivers. Positivity, the
strong Markov property at $T$, and~\eqref{growth:preparation} give
\[
 A_m(s;x)
 \ge
 \mathbb E[e^{-sT}\mathbf1_{E_0}]\,
 \frac16\,A_m^\infty(s)
 =
 \frac{A_m^\infty(s)}{12(s+3)}.
\]
This proves Theorem~\ref{growth:minorization}. The same argument applies
when other exterior neighbours are inert stubs: it only uses a real
parent's clock, its own leaf births, and nonnegative additional
immigration.

\section{An actual generation limit for the inert-parent process}

Define the uniform lower and upper sequences
\[
 u_m(s)=\inf_{x\in\mathcal X}A_m(s;x),\qquad
 v_m(s)=\sup_{x\in\mathcal X}A_m(s;x).
\]
The inert-parent state belongs to $\mathcal X$, so
\begin{equation}\label{growth:inf-comparison}
 \gamma_s A_m^\infty(s)\le u_m(s)\le A_m^\infty(s).
\end{equation}

\begin{theorem}\label{growth:inert-limit}
The limit
\[
 r_\infty(s)=\lim_{m\to\infty}A_m^\infty(s)^{1/m}
\]
exists, is finite and positive, and equals
$\lim_{m\to\infty}u_m(s)^{1/m}$. More precisely,
\begin{equation}\label{growth:sup-formula}
 \boxed{
 r_\infty(s)
 =\sup_{m\ge1}u_m(s)^{1/m}
 =\sup_{m\ge1}\bigl(\gamma_s A_m^\infty(s)\bigr)^{1/m}.
 }
\end{equation}
For every $m\ge1$,
\begin{equation}\label{growth:uniform-inert-upper}
 A_m^\infty(s)\le\gamma_s^{-1}r_\infty(s)^m.
\end{equation}
\end{theorem}
\begin{proof}
Positivity of the kernel gives
\[
 u_{m+n}\ge u_m u_n.
\]
Also, applying~\eqref{growth:comparison} to the output birth environments
of the first $m$ tagged transfers gives
\[
 A_{m+n}^\infty
 =(\mathcal K_s^m A_n)(x_\infty)
 \ge\gamma_s A_m^\infty A_n^\infty.
\]
Thus both $u_m$ and $\gamma_s A_m^\infty$ are supermultiplicative.
Their logarithms have finite linear bounds, because
$[2(s+2)]^{-m}\le A_m(s;x)\le s^{-m}$.
Fekete's lemma and~\eqref{growth:inf-comparison} prove the common limit
and both supremum formulae. The second supremum formula immediately
gives~\eqref{growth:uniform-inert-upper}.
\end{proof}

The theorem supplies a complete scalar certificate criterion for
$r_\infty(s)>1$:
\[
 r_\infty(s)>1
 \quad\Longleftrightarrow\quad
 \gamma_s A_m^\infty(s)>1
 \quad\text{for some finite }m.
\]
At rational $s$, $\gamma_s$ is rational and the existing finite-history
approximations of $A_m^\infty(s)$ can be used in exact lower certificates.
The next section identifies the inert continuous-time growth with the original exponent.

